% Answers, scoring and rubrics for assessments.tex (Level 3).
% \input by answer-key.tex -- no \documentclass here.
\kcpart{Assessments}{Assessment Answers and Rubrics}{Placement quiz, mid-level checkpoint and exit assessment.}

\begin{teachernote}[How to score]
Score with a pencil and a smile. Give the points shown; give \textbf{partial points} for a correct
idea with a small slip (for example a missing bracket). Then sit with the student for two minutes
and compare their \enquote{I can\ldots} self-rating with what you saw. The numbers are only for
you: never compare students with each other.
\end{teachernote}

\textbf{From points to the four rubric levels} (same words the students use):

\begin{center}
\renewcommand{\arraystretch}{1.3}
\begin{tabular}{@{}lcccc@{}}
\toprule
\rowcolor{kcaccent!15} Paper & Just starting & Getting there & Got it & Expert\\ \midrule
Placement quiz (20) & 0--9 & 10--15 & 16--18 & 19--20\\
Mid-level checkpoint (20) & 0--9 & 10--14 & 15--18 & 19--20\\
Exit assessment (30) & 0--12 & 13--19 & 20--25 & 26--30\\
\bottomrule
\end{tabular}
\end{center}

%% =====================================================================
\section{Placement quiz (start of week 1) --- 20 points}

\begin{answer}[Q1. Functions that give back --- 2 points]
\py|14|. \py|double(3)| gives back 6 and \py|double(4)| gives back 8.
\end{answer}

\begin{answer}[Q2. Print or return? --- 2 points]
\texttt{Hi Kai}, then \texttt{None}. (1 point for both lines, 1 for the reason.) \py|greet| prints
but never uses \py|return|, so it gives back nothing: Python's \py|None|.
\end{answer}

\begin{answer}[Q3. Lists --- 2 points]
\texttt{4 5 10}. After \py|append| there are 4 items; index 0 is the first (5), index 3 is the new 10.
\end{answer}

\begin{answer}[Q4. While loops --- 2 points]
\py|27|. \py|n| goes 1 $\rightarrow$ 3 $\rightarrow$ 9 $\rightarrow$ 27; 27 is not less than 20,
so the loop stops.
\end{answer}

\begin{answer}[Q5. Dictionaries --- 2 points]
(a) \texttt{3 \#009E73}. (b) The program crashes with a \texttt{KeyError}: there is no key
\texttt{"grass"}.
\end{answer}

\begin{answer}[Q6. Files --- 2 points]
Two lines, each \texttt{hello}: \texttt{"a"} means \emph{append} (add to the end). With
\texttt{"w"} (\emph{write}) the file is emptied first, so it would contain only one
\texttt{hello}.
\end{answer}

\begin{answer}[Q7. A pytest test --- 3 points]
One point each: a function whose name starts with \py|test_|; it calls \py|biggest| with a list;
it uses \py|assert ... == ...| with the right expected value. For example:
\begin{pycode}
def test_biggest_finds_the_largest():
    assert biggest([3, 9, 4]) == 9
\end{pycode}
Extra credit in conversation (not points): an edge case such as \py|biggest([-5, -2]) == -2|.
\end{answer}

\begin{answer}[Q8. Write a function --- 5 points]
One point each: \py|def| with the parameter; a loop over the words; the condition
\py|len(word) > 4|; a counter that goes up; \py|return| (not \py|print|) at the end.
\begin{pycode}
def count_long(words):
    count = 0
    for word in words:
        if len(word) > 4:
            count = count + 1
    return count
\end{pycode}
Also correct: \py|return len([w for w in words if len(w) > 4])|.
\end{answer}

\needspace{12\baselineskip}
\subsection{Placement decision guide}

\renewcommand{\arraystretch}{1.3}
\begin{longtable}{@{}>{\bfseries}p{2.2cm}p{\dimexpr\linewidth-2.2cm-2\tabcolsep}@{}}
\toprule
\rowcolor{kcaccent!15} Score & What to do\\ \midrule\endhead
\bottomrule\endfoot
19--20 & \textbf{Standard path + extension.} Ready for Level 3. Point them to the extension path and
stretch challenge in every lesson plan; they make good navigators for a partner who needs support.\\ \midrule
16--18 & \textbf{Standard path.} Ready. Glance at which question was missed and mention it when it
comes up again.\\ \midrule
10--15 & \textbf{Standard path with support.} Start Level 3, but use the lesson plans' support paths
in weeks 1--3, and set the refresher below as extra homework for each missed topic. Re-ask the
missed questions orally in week 3.\\ \midrule
0--9 & \textbf{Refresher first.} Before week 1 (or using one buffer session), redo the Level 2
homework listed below together. Pair this student with you or a confident, patient partner in
weeks 1--5. Then let them retake only the questions they missed.
If two or more students score below 10, consider repeating the Level 2 final project first.\\
\end{longtable}

\textbf{Refresher by question} (Level 2 materials):
Q1, Q2, Q8 $\rightarrow$ week 2 (functions that give back);
Q3 $\rightarrow$ week 3 (lists);
Q4 $\rightarrow$ week 4 (while loops);
Q5 $\rightarrow$ week 5 (dictionaries);
Q6 $\rightarrow$ week 7 (files);
Q7 $\rightarrow$ week 8 (testing with pytest).

\begin{tip}[Self-rating mismatch]
If a student checks \enquote{Got it} or \enquote{Expert} but missed the questions on that topic, do
not argue: say \enquote{Let's check together} and do one similar example side by side. If a
student rates themself low but scored well, tell them so, specifically.
\end{tip}

%% =====================================================================
\section{Mid-level checkpoint (week 6) --- 20 points}

\begin{answer}[Q1. Read a class --- 2 points]
\texttt{Mo had fun!} then \texttt{1}. Energy starts at 5 and \py|play()| runs twice: 5 $-$ 2 $-$ 2 = 1.
\end{answer}

\begin{answer}[Q2. Families of objects --- 4 points]
(a) \texttt{7 8} and \texttt{Omar says hi!} (2 points). Zoe loses 3 once; Omar loses 1 twice.
(b) \py|hit| (1 point). (c) \py|shout| is \textbf{inherited} from \py|Character|: \py|Knight|
does not define it, so Python looks in the parent class (1 point).
\end{answer}

\begin{answer}[Q3. Fix the bug --- 2 points]
Line 3: \py|count = 0| makes an ordinary variable that disappears when \py|__init__| ends. It must
be an attribute of the object: \py|self.count = 0|.
\end{answer}

\begin{answer}[Q4. Trace a binary search --- 5 points]
\begin{center}
\renewcommand{\arraystretch}{1.3}
\begin{tabular}{|c|c|c|c|l|}
\hline
\textbf{Step} & \textbf{low} & \textbf{high} & \textbf{middle} & \textbf{item at middle}\\ \hline
1 & 0 & 9 & 4 & 16: too small, so low = 5\\ \hline
2 & 5 & 9 & 7 & 56: too big, so high = 6\\ \hline
3 & 5 & 6 & 5 & 23: found!\\ \hline
\end{tabular}
\end{center}
(a) 1 point per correct row. (b) Linear search checks \textbf{6} items (indexes 0--5) (1 point).
(c) About \textbf{10} (1 point): each check halves the list, and halving 1000 about ten times gets
down to 1. Accept 9--11.
\end{answer}

\begin{answer}[Q5. Sorting --- 3 points]
Pass 1: \texttt{[1, 3, 9, 7]}; pass 2: \texttt{[1, 3, 9, 7]} (3 is already in place);
pass 3: \texttt{[1, 3, 7, 9]} (2 points; 1 point if one pass is wrong). The program prints
\texttt{Leo}: sorted by score, biggest first (1 point).
\end{answer}

\begin{answer}[Q6. Stacks and queues --- 2 points]
\texttt{C} then \texttt{A}: a stack gives back the \emph{last} item added, a queue the
\emph{first} (1 point). Undo button: S; players waiting: Q; pile of plates: S (1 point).
\end{answer}

\begin{answer}[Q7. Explain --- 2 points]
Each check of binary search throws away half of what is left, so the number of checks grows very
slowly; linear search checks items one by one (1 point). The list must be \textbf{sorted} (1 point).
\end{answer}

\begin{rubric}[What to do after the checkpoint]
\textbf{Mostly Q1--Q3 missed (classes):} in the week 6 review, rebuild the week 2 pet class
together, then add a child class. \textbf{Mostly Q4--Q5 missed (algorithms):} schedule buffer B1
(Algorithm Olympics) before week 7. \textbf{Q6 missed:} replay the plate-stack and lunch-queue
unplugged activity for 5 minutes. \textbf{19--20:} ready for the second half; point them to each lesson's extension path and
stretch challenge from week 7 on.
\end{rubric}

%% =====================================================================
\section{Exit assessment (week 12) --- 30 points}

\begin{answer}[Q1. Recursion --- 2 points]
\texttt{3}, \texttt{2}, \texttt{1}, \texttt{Blast off!} on four lines.
\end{answer}

\begin{answer}[Q2. Recursion that returns --- 2 points]
\py|10|: \py|total(4)| = 4 + \py|total(3)| = 4 + 3 + \py|total(2)| = 4 + 3 + 2 + \py|total(1)| =
4 + 3 + 2 + 1. One point for the answer, one for showing the chain of calls.
\end{answer}

\begin{answer}[Q3. Grids --- 3 points]
(a) \texttt{F G}. (b) \texttt{grid[2][1]}: the third row and the second column, because counting starts at 0. (c)
\texttt{["A", "\#", "C"]}.
\end{answer}

\begin{answer}[Q4. A recursion bug --- 4 points]
There is no \textbf{base case}, so \py|stars| calls itself forever (4, 3, 2, 1, 0, $-1$, \ldots)
until Python stops it with \texttt{RecursionError} (2 points). A fix (2 points):
\begin{pycode}
def stars(n):
    if n == 0:          # base case: stop here
        return
    print("*" * n)
    stars(n - 1)        # a smaller problem each time
\end{pycode}
Also correct: \py|if n > 0:| around the two lines.
\end{answer}

\begin{answer}[Q5. A data bug --- 3 points]
Everything read from a CSV is text, so line 6 becomes \py|total = total + int(row["points"])|
(2 points). One sensible guard (1 point): skip empty cells
(\py|if row["points"].strip() == "": continue|), strip spaces, check \py|.isdigit()|, or open
the file with \py|encoding="utf-8-sig"| if it came from a spreadsheet app.
\end{answer}

\begin{answer}[Q6. Nearest neighbor --- 4 points]
The paper uses this simpler distance (add the two gaps) instead of the ruler distance from
week 9, so no square roots are needed. Distances to the mystery monster (size 4, spikes 1): A = 5 + 3 = \textbf{8}; B = 4 + 4 = \textbf{8};
C = 2 + 1 = \textbf{3}; D = 1 + 0 = \textbf{1} (1 point per two correct distances, so 2 points).
Prediction: \textbf{slime}, because the nearest example, D, is a slime (2 points: label + reason).
\end{answer}

\begin{answer}[Q7. When AI gets it wrong --- 3 points]
(a) It may say it is not an apple, or give it the wrong label, because it has never seen a green
apple. (b) \textbf{Bias} / biased or unrepresentative training data. (c) Train it on many
different apples (green, yellow, red, big, small) and test it on examples it has not seen.
\end{answer}

\begin{answer}[Q8. Using AI safely --- 3 points]
(a) Any two of: real name, home address, school, phone number, passwords, photos of themselves or
others, other people's private information (1 point). (b) It only counts which word followed which
in its training text and picks the next word by chance; it has no idea what the words mean and
cannot check facts (2 points). Chatbots are far more advanced but can also sound sure and be wrong.
\end{answer}

\begin{answer}[Q9. Design a game --- 4 points]
One point each. (a) A \textbf{2D grid} (list of lists, \texttt{grid[row][col]}). (b)
\textbf{Breadth-first search} (BFS), which needs a \textbf{queue} (\texttt{collections.deque}).
BFS is met in buffer B1 and in the project (Maze Muncher's Chaser), not in a core lesson, so also
give the point for a clear description such as \enquote{spread out one square at a time, like the
paint bucket, remembering where each square came from, using a queue}.
(c) Two sensible classes with an attribute and a method each, for example \texttt{Player}
(\texttt{pos}; \texttt{move()}) and \texttt{Monster} (\texttt{pos}; \texttt{chase()}). The fourth point is for a good
reason: both could inherit from a \texttt{Character} base class because each \emph{is a}
character; a \texttt{Game} or \texttt{Map} class \emph{has} the grid, the player and the monster
(composition). Accept any well-argued design.
\end{answer}

\begin{answer}[Q10. Your design decision --- 2 points]
One point for naming a real decision and an alternative; one point for a reason (simpler, easier
to test, faster, fairer to players, easier to add levels\ldots). Example: \enquote{I kept the game
rules in a separate file with no turtle, so I could test them with pytest. I could have put
everything in one file, but then my tests would open windows.}
\end{answer}

\begin{rubric}[Exit assessment by Level 3 goal]
Use this to fill in the progress tracker's final column.
\begin{center}
\renewcommand{\arraystretch}{1.25}
\begin{tabularx}{\linewidth}{@{}>{\raggedright\arraybackslash}p{4.6cm}>{\raggedright\arraybackslash}p{2.3cm}>{\raggedright\arraybackslash}X@{}}
\toprule
Goal & Questions & Look for\\ \midrule
L3.G1 Classes \& design & Q9, Q10 & sensible classes; inheritance vs composition explained\\
L3.G2 Algorithms \& recursion & Q1, Q2, Q4 & base case; smaller problem each call\\
L3.G3 Data structures & Q3, Q9 & grid indexing; queue for BFS\\
L3.G5 Iterative project & Q10 & a real decision with a reason\\
L3.G6 Data & Q5 & text vs numbers; cleaning\\
L3.G7 AI & Q6, Q7, Q8 & distance, nearest label, bias\\
L3.G8 Test \& debug & Q4, Q5 & reads the error, fixes the cause\\
L3.G9 Ethics \& privacy & Q7, Q8 & fairness of data; what stays private\\
\bottomrule
\end{tabularx}
\end{center}
L3.G4 (git) is assessed through the project's commit history, not on paper.
\end{rubric}

\begin{teachernote}[Graduation is not a score]
Every student who builds and presents a final project graduates from Level 3 and receives the
certificate. Use the exit score and project rubric to write one specific sentence of praise and
one \enquote{next step} on the back of the certificate or in the progress tracker.
\end{teachernote}
